% --- Archon physics formalization source begin ---
% archon:physics
% archon:covers PhyXMiniProblems/problem_phyx_mini_0809.lean
% archon:source-report reports/phyx_mini/problem_phyx_mini_0809.source.json

\chapter{Physics problem phyx\_mini\_0809}
\label{ch:PhyXMiniProblems_problem_phyx_mini_0809}

\paragraph{Problem source.}
\#\# Physical scenario

At a family picnic you are appointed to push your obnoxious cousin Throckmorton in a swing as shown in figure His weight is \textbackslash{}( w \textbackslash{}), the length of the chains is \textbackslash{}( R \textbackslash{}), and you push Throcky until the chains make an angle \textbackslash{}( \textbackslash{}theta\_0 \textbackslash{}) with the vertical. To do this, you exert a varying horizontal force \textbackslash{}( \textbackslash{}vec\{F\} \textbackslash{}) that starts at zero and gradually increases just enough that Throcky and the swing move very slowly and remain very nearly in equilibrium throughout the process.

\#\# Question

What is the work you do by exerting force \textbackslash{}( \textbackslash{}vec\{F\} \textbackslash{})? (Ignore the weight of the chains and seat.)

\#\# Answer choices

- A: A: \textbackslash{}( 2wR(1 - \textbackslash{}cos \textbackslash{}theta\_0) \textbackslash{})

- B: B: \textbackslash{}( wR(1 - \textbackslash{}cos \textbackslash{}theta\_0) \textbackslash{})

- C: C: \textbackslash{}( wR(1 - \textbackslash{}sin \textbackslash{}theta\_0) \textbackslash{})

- D: D: \textbackslash{}( 2wR(1 - \textbackslash{}sin \textbackslash{}theta\_0) \textbackslash{})

\#\# Figure caption (auxiliary; use the image as primary evidence)

The image depicts a person sitting on a swing, which is hanging from a chain. The swing is displaced from its vertical position, creating an angle labeled as \textbackslash{}( \textbackslash{}theta \textbackslash{}) with the dashed lines, which illustrate the swing’s path.

Key elements in the image:
- **Person**: Seated on the swing, holding the chains. The person is wearing a cap with a "B" on it.
- **Swing Path**: Illustrated with a blue dashed line in an arc, denoting the path or trajectory of the swing.
- **Forces and Vectors**:
  - A red arrow labeled \textbackslash{}( \textbackslash{}vec\{F\} \textbackslash{}) represents a force applied horizontally.
  - A blue arrow marked \textbackslash{}( d\textbackslash{}vec\{l\} \textbackslash{}) represents an infinitesimal path element along the arc.
- **Angle \textbackslash{}( \textbackslash{}theta \textbackslash{})**: This angle is shown between the vertical dashed line and the swing, as well as between the person's current position and the lowest point of the arc.
- **Radius \textbackslash{}( R \textbackslash{})**: The vertical dashed line labeled with \textbackslash{}( R \textbackslash{}) shows the length from the pivot point to the lowest point.
- **Arc Length \textbackslash{}( s \textbackslash{})**: Represented by the distance along the arc from the resting vertical position to the current position.

Overall, the image illustrates the physics of a pendulum, showing the forces, displacement, and path of a swinging motion.

\#\# Dataset metadata

Work and Energy; Physical Model Grounding Reasoning, Numerical Reasoning

\paragraph{Recorded answer/context.}
B: B: \textbackslash{}( wR(1 - \textbackslash{}cos \textbackslash{}theta\_0) \textbackslash{})

\paragraph{Figure/image path.}
/root/proposal\_for\_physic/science-mango/phyx\_mini\_run/phyx\_data/test\_image/809.png

\paragraph{Formalization target.}
create a compiling Lean file with sorry bodies at `PhyXMiniProblems/problem\_phyx\_mini\_0809.lean`. The Lean declarations must preserve the physical quantities, dimensions or dimensional roles, figure labels, governing-law hypotheses, and final relation expressed by this problem.
Use Mathlib/Physlib names found through LeanExplore where available. If a physics API is missing, introduce faithful local abstractions rather than scalar placeholder aliases.

\begin{theorem}[Physics formalization target]
\label{thm:physics:phyx_mini_0809:target}
\lean{PhyXMiniProblems.ProblemPhyXMini0809.problem_phyx_mini_0809}
\uses{def:physics:phyx-mini-0809:phyxminiproblems-problemphyxmini0809-lengthreadout, def:physics:phyx-mini-0809:phyxminiproblems-problemphyxmini0809-forcemagnitudereadout, def:physics:phyx-mini-0809:phyxminiproblems-problemphyxmini0809-workreadout, def:physics:phyx-mini-0809:phyxminiproblems-problemphyxmini0809-lengthinmeters, def:physics:phyx-mini-0809:phyxminiproblems-problemphyxmini0809-forcemagnitudeinnewtons, def:physics:phyx-mini-0809:phyxminiproblems-problemphyxmini0809-workinjoules, def:physics:phyx-mini-0809:phyxminiproblems-problemphyxmini0809-quasistaticswingsetup, def:physics:phyx-mini-0809:phyxminiproblems-problemphyxmini0809-matchesproblemandprimaryfigure, def:physics:phyx-mini-0809:phyxminiproblems-problemphyxmini0809-hasphysicalswingparameters, def:physics:phyx-mini-0809:phyxminiproblems-problemphyxmini0809-satisfiescircularswinggeometry, def:physics:phyx-mini-0809:phyxminiproblems-problemphyxmini0809-satisfiesquasistatictangentialequilibrium, def:physics:phyx-mini-0809:phyxminiproblems-problemphyxmini0809-satisfiesappliedforceworklaw, def:physics:phyx-mini-0809:phyxminiproblems-problemphyxmini0809-matchesdisplayedanswer}
The assigned autoformalize agent should translate this physics problem into Lean declarations in the covered file, with theorem and lemma proof bodies written as `by sorry`.
\end{theorem}
\begin{proof}
This is an autoformalization task, not a proof task. Produce statements that can later be proved by the physics prover without weakening the physical model.
\end{proof}

% --- Archon declaration topology begin ---
\section{Declaration topology}
\begin{definition}[force Dimension]
  \label{def:physics:phyx-mini-0809:phyxminiproblems-problemphyxmini0809-forcedimension}
  \lean{PhyXMiniProblems.ProblemPhyXMini0809.forceDimension}
  The physical dimension of a force magnitude, M L T⁻²
\end{definition}
\begin{proof}
  This is the declared model definition of force Dimension.
\end{proof}
\begin{definition}[Length Quantity]
  \label{def:physics:phyx-mini-0809:phyxminiproblems-problemphyxmini0809-lengthquantity}
  \lean{PhyXMiniProblems.ProblemPhyXMini0809.LengthQuantity}
  A nonnegative, unit-independent physical length
\end{definition}
\begin{proof}
  This is the declared model definition of Length Quantity.
\end{proof}
\begin{definition}[Force Magnitude]
  \label{def:physics:phyx-mini-0809:phyxminiproblems-problemphyxmini0809-forcemagnitude}
  \lean{PhyXMiniProblems.ProblemPhyXMini0809.ForceMagnitude}
  \uses{def:physics:phyx-mini-0809:phyxminiproblems-problemphyxmini0809-forcedimension}
  A nonnegative, unit-independent physical force magnitude
\end{definition}
\begin{proof}
  This is the declared model definition of Force Magnitude.
\end{proof}
\begin{definition}[Work Quantity]
  \label{def:physics:phyx-mini-0809:phyxminiproblems-problemphyxmini0809-workquantity}
  \lean{PhyXMiniProblems.ProblemPhyXMini0809.WorkQuantity}
  Signed physical work, carrying Physlib's energy dimension
\end{definition}
\begin{proof}
  This is the declared model definition of Work Quantity.
\end{proof}
\begin{definition}[length Readout]
  \label{def:physics:phyx-mini-0809:phyxminiproblems-problemphyxmini0809-lengthreadout}
  \lean{PhyXMiniProblems.ProblemPhyXMini0809.lengthReadout}
  \uses{def:physics:phyx-mini-0809:phyxminiproblems-problemphyxmini0809-lengthquantity}
  Scalar readout of a physical length in a coherent choice of base units
\end{definition}
\begin{proof}
  This is the declared model definition of length Readout.
\end{proof}
\begin{definition}[force Magnitude Readout]
  \label{def:physics:phyx-mini-0809:phyxminiproblems-problemphyxmini0809-forcemagnitudereadout}
  \lean{PhyXMiniProblems.ProblemPhyXMini0809.forceMagnitudeReadout}
  \uses{def:physics:phyx-mini-0809:phyxminiproblems-problemphyxmini0809-forcemagnitude}
  Scalar readout of a force magnitude in a coherent choice of base units
\end{definition}
\begin{proof}
  This is the declared model definition of force Magnitude Readout.
\end{proof}
\begin{definition}[work Readout]
  \label{def:physics:phyx-mini-0809:phyxminiproblems-problemphyxmini0809-workreadout}
  \lean{PhyXMiniProblems.ProblemPhyXMini0809.workReadout}
  \uses{def:physics:phyx-mini-0809:phyxminiproblems-problemphyxmini0809-workquantity}
  Signed scalar readout of work in a coherent choice of base units
\end{definition}
\begin{proof}
  This is the declared model definition of work Readout.
\end{proof}
\begin{definition}[length In Meters]
  \label{def:physics:phyx-mini-0809:phyxminiproblems-problemphyxmini0809-lengthinmeters}
  \lean{PhyXMiniProblems.ProblemPhyXMini0809.lengthInMeters}
  \uses{def:physics:phyx-mini-0809:phyxminiproblems-problemphyxmini0809-lengthquantity, def:physics:phyx-mini-0809:phyxminiproblems-problemphyxmini0809-lengthreadout}
  Metre readout of a physical length in coherent SI units
\end{definition}
\begin{proof}
  This is the declared model definition of length In Meters.
\end{proof}
\begin{definition}[force Magnitude In Newtons]
  \label{def:physics:phyx-mini-0809:phyxminiproblems-problemphyxmini0809-forcemagnitudeinnewtons}
  \lean{PhyXMiniProblems.ProblemPhyXMini0809.forceMagnitudeInNewtons}
  \uses{def:physics:phyx-mini-0809:phyxminiproblems-problemphyxmini0809-forcemagnitude, def:physics:phyx-mini-0809:phyxminiproblems-problemphyxmini0809-forcemagnitudereadout}
  Newton readout of a force magnitude in coherent SI units
\end{definition}
\begin{proof}
  This is the declared model definition of force Magnitude In Newtons.
\end{proof}
\begin{definition}[work In Joules]
  \label{def:physics:phyx-mini-0809:phyxminiproblems-problemphyxmini0809-workinjoules}
  \lean{PhyXMiniProblems.ProblemPhyXMini0809.workInJoules}
  \uses{def:physics:phyx-mini-0809:phyxminiproblems-problemphyxmini0809-workquantity, def:physics:phyx-mini-0809:phyxminiproblems-problemphyxmini0809-workreadout}
  Joule readout of signed physical work
\end{definition}
\begin{proof}
  This is the declared model definition of work In Joules.
\end{proof}
\begin{definition}[Swing Weight Model]
  \label{def:physics:phyx-mini-0809:phyxminiproblems-problemphyxmini0809-swingweightmodel}
  \lean{PhyXMiniProblems.ProblemPhyXMini0809.SwingWeightModel}
  Idealized weight models relevant to the problem's instruction
\end{definition}
\begin{proof}
  This is the declared model definition of Swing Weight Model.
\end{proof}
\begin{definition}[Swing Process Regime]
  \label{def:physics:phyx-mini-0809:phyxminiproblems-problemphyxmini0809-swingprocessregime}
  \lean{PhyXMiniProblems.ProblemPhyXMini0809.SwingProcessRegime}
  The slow process stipulated in the prose problem
\end{definition}
\begin{proof}
  This is the declared model definition of Swing Process Regime.
\end{proof}
\begin{definition}[Figure Element]
  \label{def:physics:phyx-mini-0809:phyxminiproblems-problemphyxmini0809-figureelement}
  \lean{PhyXMiniProblems.ProblemPhyXMini0809.FigureElement}
  Named objects and curves visible in the supplied bitmap
\end{definition}
\begin{proof}
  This is the declared model definition of Figure Element.
\end{proof}
\begin{definition}[Figure Vector Label]
  \label{def:physics:phyx-mini-0809:phyxminiproblems-problemphyxmini0809-figurevectorlabel}
  \lean{PhyXMiniProblems.ProblemPhyXMini0809.FigureVectorLabel}
  The two vector arrows explicitly labeled in the supplied bitmap
\end{definition}
\begin{proof}
  This is the declared model definition of Figure Vector Label.
\end{proof}
\begin{definition}[Figure Vector Direction]
  \label{def:physics:phyx-mini-0809:phyxminiproblems-problemphyxmini0809-figurevectordirection}
  \lean{PhyXMiniProblems.ProblemPhyXMini0809.FigureVectorDirection}
  Qualitative directions needed to transcribe the two pictured arrows
\end{definition}
\begin{proof}
  This is the declared model definition of Figure Vector Direction.
\end{proof}
\begin{definition}[Swing Figure]
  \label{def:physics:phyx-mini-0809:phyxminiproblems-problemphyxmini0809-swingfigure}
  \lean{PhyXMiniProblems.ProblemPhyXMini0809.SwingFigure}
  \uses{def:physics:phyx-mini-0809:phyxminiproblems-problemphyxmini0809-figureelement, def:physics:phyx-mini-0809:phyxminiproblems-problemphyxmini0809-figurevectorlabel, def:physics:phyx-mini-0809:phyxminiproblems-problemphyxmini0809-figurevectordirection}
  Direct qualitative and textual evidence from the primary figure
\end{definition}
\begin{proof}
  This is the declared model definition of Swing Figure.
\end{proof}
\begin{definition}[Quasistatic Swing Setup]
  \label{def:physics:phyx-mini-0809:phyxminiproblems-problemphyxmini0809-quasistaticswingsetup}
  \lean{PhyXMiniProblems.ProblemPhyXMini0809.QuasistaticSwingSetup}
  \uses{def:physics:phyx-mini-0809:phyxminiproblems-problemphyxmini0809-lengthquantity, def:physics:phyx-mini-0809:phyxminiproblems-problemphyxmini0809-forcemagnitude, def:physics:phyx-mini-0809:phyxminiproblems-problemphyxmini0809-workquantity, def:physics:phyx-mini-0809:phyxminiproblems-problemphyxmini0809-swingweightmodel, def:physics:phyx-mini-0809:phyxminiproblems-problemphyxmini0809-swingprocessregime, def:physics:phyx-mini-0809:phyxminiproblems-problemphyxmini0809-swingfigure}
  The independent physical quantities and path data. The force magnitude is a function of the current angular position because the push is varied during the motion. Neither workDoneByPush nor any field is defined from an answer choice or from the requested closed formula
\end{definition}
\begin{proof}
  This is the declared model definition of Quasistatic Swing Setup.
\end{proof}
\begin{definition}[Matches Problem And Primary Figure]
  \label{def:physics:phyx-mini-0809:phyxminiproblems-problemphyxmini0809-matchesproblemandprimaryfigure}
  \lean{PhyXMiniProblems.ProblemPhyXMini0809.MatchesProblemAndPrimaryFigure}
  \uses{def:physics:phyx-mini-0809:phyxminiproblems-problemphyxmini0809-forcemagnitudeinnewtons, def:physics:phyx-mini-0809:phyxminiproblems-problemphyxmini0809-figureelement, def:physics:phyx-mini-0809:phyxminiproblems-problemphyxmini0809-quasistaticswingsetup}
  Problem-statement and primary-image readouts. The swing starts at the lowest point (θ = 0), the horizontal push starts from zero, only the rider's weight is retained, and all labels and directions are transcribed from the bitmap. No value of the requested work occurs here
\end{definition}
\begin{proof}
  This is the declared model definition of Matches Problem And Primary Figure.
\end{proof}
\begin{definition}[Has Physical Swing Parameters]
  \label{def:physics:phyx-mini-0809:phyxminiproblems-problemphyxmini0809-hasphysicalswingparameters}
  \lean{PhyXMiniProblems.ProblemPhyXMini0809.HasPhysicalSwingParameters}
  \uses{def:physics:phyx-mini-0809:phyxminiproblems-problemphyxmini0809-lengthinmeters, def:physics:phyx-mini-0809:phyxminiproblems-problemphyxmini0809-forcemagnitudeinnewtons, def:physics:phyx-mini-0809:phyxminiproblems-problemphyxmini0809-quasistaticswingsetup}
  Positivity and the acute-angle branch depicted in the figure. This is the branch on which a finite horizontal force can maintain static equilibrium
\end{definition}
\begin{proof}
  This is the declared model definition of Has Physical Swing Parameters.
\end{proof}
\begin{definition}[Satisfies Circular Swing Geometry]
  \label{def:physics:phyx-mini-0809:phyxminiproblems-problemphyxmini0809-satisfiescircularswinggeometry}
  \lean{PhyXMiniProblems.ProblemPhyXMini0809.SatisfiesCircularSwingGeometry}
  \uses{def:physics:phyx-mini-0809:phyxminiproblems-problemphyxmini0809-lengthreadout, def:physics:phyx-mini-0809:phyxminiproblems-problemphyxmini0809-quasistaticswingsetup}
  Circular-arc geometry from the primary figure: s = R (θ - θᵢ), and the forward tangent is inclined by θ above the horizontal when the chain is inclined by θ from the vertical. Both statements hold throughout the depicted push
\end{definition}
\begin{proof}
  This is the declared model definition of Satisfies Circular Swing Geometry.
\end{proof}
\begin{definition}[Satisfies Quasistatic Tangential Equilibrium]
  \label{def:physics:phyx-mini-0809:phyxminiproblems-problemphyxmini0809-satisfiesquasistatictangentialequilibrium}
  \lean{PhyXMiniProblems.ProblemPhyXMini0809.SatisfiesQuasistaticTangentialEquilibrium}
  \uses{def:physics:phyx-mini-0809:phyxminiproblems-problemphyxmini0809-forcemagnitudereadout, def:physics:phyx-mini-0809:phyxminiproblems-problemphyxmini0809-quasistaticswingsetup}
  Tangential static equilibrium. Chain tension has no tangential component, so the component of the horizontal push along the forward tangent balances the opposing tangential component of the rider's weight. This local law contains no integrated work or endpoint work formula
\end{definition}
\begin{proof}
  This is the declared model definition of Satisfies Quasistatic Tangential Equilibrium.
\end{proof}
\begin{definition}[Satisfies Applied Force Work Law]
  \label{def:physics:phyx-mini-0809:phyxminiproblems-problemphyxmini0809-satisfiesappliedforceworklaw}
  \lean{PhyXMiniProblems.ProblemPhyXMini0809.SatisfiesAppliedForceWorkLaw}
  \uses{def:physics:phyx-mini-0809:phyxminiproblems-problemphyxmini0809-lengthreadout, def:physics:phyx-mini-0809:phyxminiproblems-problemphyxmini0809-forcemagnitudereadout, def:physics:phyx-mini-0809:phyxminiproblems-problemphyxmini0809-workreadout, def:physics:phyx-mini-0809:phyxminiproblems-problemphyxmini0809-quasistaticswingsetup}
  Work done by the applied force is its line integral along the circular path. With θ as path parameter, |d⃗l| = R dθ; the cosine is the dot-product projection of the horizontal force onto the forward tangent. The law is stated in every coherent unit system and does not evaluate the integral
\end{definition}
\begin{proof}
  This is the declared model definition of Satisfies Applied Force Work Law.
\end{proof}
\begin{definition}[Answer Choice]
  \label{def:physics:phyx-mini-0809:phyxminiproblems-problemphyxmini0809-answerchoice}
  \lean{PhyXMiniProblems.ProblemPhyXMini0809.AnswerChoice}
  Labels of the four answer choices printed with the problem
\end{definition}
\begin{proof}
  This is the declared model definition of Answer Choice.
\end{proof}
\begin{definition}[displayed Work In Joules]
  \label{def:physics:phyx-mini-0809:phyxminiproblems-problemphyxmini0809-displayedworkinjoules}
  \lean{PhyXMiniProblems.ProblemPhyXMini0809.displayedWorkInJoules}
  \uses{def:physics:phyx-mini-0809:phyxminiproblems-problemphyxmini0809-lengthinmeters, def:physics:phyx-mini-0809:phyxminiproblems-problemphyxmini0809-forcemagnitudeinnewtons, def:physics:phyx-mini-0809:phyxminiproblems-problemphyxmini0809-quasistaticswingsetup, def:physics:phyx-mini-0809:phyxminiproblems-problemphyxmini0809-answerchoice}
  The symbolic work expression printed beside each answer label, evaluated using coherent-SI newton, metre, and joule readouts
\end{definition}
\begin{proof}
  This is the declared model definition of displayed Work In Joules.
\end{proof}
\begin{definition}[Matches Displayed Answer]
  \label{def:physics:phyx-mini-0809:phyxminiproblems-problemphyxmini0809-matchesdisplayedanswer}
  \lean{PhyXMiniProblems.ProblemPhyXMini0809.MatchesDisplayedAnswer}
  \uses{def:physics:phyx-mini-0809:phyxminiproblems-problemphyxmini0809-workinjoules, def:physics:phyx-mini-0809:phyxminiproblems-problemphyxmini0809-quasistaticswingsetup, def:physics:phyx-mini-0809:phyxminiproblems-problemphyxmini0809-answerchoice, def:physics:phyx-mini-0809:phyxminiproblems-problemphyxmini0809-displayedworkinjoules}
  The physical work agrees with the symbolic expression of a displayed choice
\end{definition}
\begin{proof}
  This is the declared model definition of Matches Displayed Answer.
\end{proof}
% --- Archon declaration topology end ---
% --- Archon physics formalization source end ---
